The strongest result supports a content-allocation interpretation: when many details are fixed independently of the secret, fewer free choices remain for secret-dependent associations to influence. The weak filler effect makes pure context length an insufficient explanation in this setup. However, the experiment cannot distinguish commitment to a plan from ordinary semantic priming by the supplied events, and reduced discrimination is not an information-theoretic secrecy guarantee.

The plot findings narrow the interpretation further. Already providing a premise creates a much more constrained task than open-ended word-secret writing, and the generic opening plan adds less specificity than the concrete word outlines. A weaker plot effect is compatible with limited remaining content freedom and a weaker baseline signal; it does not demonstrate a fundamental distinction between words and twists. The residual results likewise separate detectable context dependence from successful causal erasure. A residual direction can contain word-related information while a single projection intervention fails to eliminate other information paths, or even disrupt instructions that prevent literal disclosure.
